\documentclass{article} %210 mm × 297 mm
\usepackage[utf8]{inputenc}
\usepackage[a4paper,left=1.5cm, right=1.5cm, top=2cm, bottom=2cm]{geometry}
\usepackage{graphicx}
%\usepackage{blindtext}
\usepackage{multirow} 
\usepackage{mhchem}
\usepackage{url}
\usepackage{subfigure}
\usepackage{amsfonts,latexsym} % para tener disponibilidad de diversos simbolos
\usepackage{enumerate}
%\usepackage{booktabs}
\usepackage{float}
%\usepackage{threeparttable}
\usepackage{array,colortbl}
%\usepackage{ifpdf}
%\usepackage{rotating}
\usepackage{cite}
\usepackage{stfloats}
\usepackage{url}
\usepackage{listings}
\usepackage{fancyhdr}

\usepackage{color}
%\usepackage{hyperref}
%\usepackage{wrapfig}
\usepackage{array}
%\usepackage{multirow}
%\usepackage{adjustbox}
%\usepackage{nccmath}
\usepackage[dvipsnames]{xcolor}


\newcommand{\titulo}{Präsenzübungsblatt 1; MfN II.}
\newcommand{\nombre}{Liliana Sanfilippo}
\newcommand{\FCE}{\textbf{SoSe 2025}}

 

\usepackage{fancyhdr}
\pagestyle{fancy}
\fancyhead{}
\fancyfoot{}
\fancyhead[LO]{\small\nombre}
\fancyhead[RE]{\small\titulo}
\fancyhead[LO]{\small\FCE}
\fancyfoot[RO,LE] {\thepage}

\setcounter{page}{1} %Número de la página, la editorial lo cambiará



\usepackage[onehalfspacing]{setspace}
\begin{document}


\begin{tabular}{p{12cm}}
{{\Large\bf\titulo}}\\
\vspace{0.2cm}\\
 {\large\nombre}\\
\end{tabular}
\vspace{0.2cm}\\
\section{Aufgabe 1.1}
\textcolor{Blue}{Was ist überhaupt $\Phi$ hier? Eine lineare Abbildung $\Phi: V \rightarrow V$. \\
Was ist $V$ hier? ein Vektorraum über einem Körper $K$. \\
Wie hängt eine Matrix damit zusammen? Man kann lineare Abbildungen durch Matrizen darstellen. \\
Was passiert, wenn man lineare Abbildungen mit sich selber multipliziert ($\Phi^n$)? Multiplikation linearer Abbildungen ist eigentlich eine hintereinander her Ausführung der Funktion, also eine Verkettung $\Phi(\Phi(\cdots))$. Weil man lineare Abbildungen durch Matrizen darstellen kann, entspricht das einer Matrixmultiplikation einer Matrix mit sich selber. \\
Was für Matrizen muss man Multiplizieren, damit 0 raus kommt? \\
Wie kann es überhaupt sein, dass $A \cdot v = 0 \cdot v $? \\
Was bedeutet es, wenn der Eigenwert 0 ist? 
} \\
Sei $a$ EW von $\Phi$.
$\Phi(v) = a \cdot v$
Wenn $\Phi(v) = 0$, ist auch $a^n \cdot v = 0$, also muss $a^n = 0$ und auch $a = 0$
\section{Aufgabe 1.2}
$F^2 + F$ hat EW -1 \\
Sei $v \neq 0$ EV zu -1 \\
\begin{align*}
    F^2(v) + F(v) & = -v  \\
    \Rightarrow F^2(v) + F(v) + v & = 0 \\
    \overset{F}{\Rightarrow} F^3(v) + \underbrace{F^2(v) + F(v)}_{=-v} & = 0 \\
    \Rightarrow F^3(v) - v & = 0 \\
    \Rightarrow F^3(v) & = v \\
    \Rightarrow F \text{ hat EW } 1
\end{align*}


\section{Aufgabe 1.3}
\subsubsection{A.}
\begin{align*}
    A & = \begin{pmatrix}2 & -1 \\ 1 & 4 \end{pmatrix} \\ \\
    det(A - \lambda \cdot E) & = (2-\lambda) (4-\lambda) +1 \\
    & = \lambda^2 - 6 \lambda +9 \\ \\
    \lambda_{1|2} & = 3 \overset{+}{-} \sqrt{0} \\
    \Rightarrow \lambda_{1|2} & = 3
\end{align*}
EV zu EW $\lambda = 3$
\begin{align*}
    \begin{pmatrix}
    2-3 & -1 \\ 1 & 4-3
    \end{pmatrix} \rightsquigarrow \begin{pmatrix}
    -1 & -1 \\ 1 & 1 
    \end{pmatrix} \rightsquigarrow \begin{pmatrix}
    -1 & -1  \\ 0 & 0 
    \end{pmatrix} \Rightarrow x_2 beliebig  \\ \\
    x_1 = -x_2 \\ 
    \mathcal{C} = \Biggl \{ \binom{-1}{1} \Biggr \}
\end{align*}
\subsubsection{B.}
\begin{align*}
    B & = \begin{pmatrix}
    1 & -3i \\ i & -1
    \end{pmatrix} \\ \\
    det(B - \lambda E) & = (1-\lambda) (-1-\lambda) - (-3i^2) \\
    & = \lambda^2 -4 \\ \\
    \lambda^2 & = 4 \overset{\sqrt{}}{\Rightarrow} \lambda_1 = 2, \quad \lambda_2 = -2 \\ 
\end{align*}
EV zu EW $\lambda = -2$
\begin{align*}
    \begin{pmatrix}
    2 & -3i \\ i & 1 
    \end{pmatrix} \overset{z2 + z1 \cdot \frac{-i}{3}}{\rightsquigarrow }\begin{pmatrix}
    3 & -3i \\ 0 & 0 
    \end{pmatrix} \Rightarrow x_2 beliebig \\ \\
    x_1 & = i x_2 \\
     \mathcal{C} = \Biggl \{ \binom{i}{1} \Biggr \}
\end{align*}
EV zu EW $\lambda = 2$
\begin{align*}
    \begin{pmatrix}
    -1 & -3i \\ i & -3 
    \end{pmatrix}  \overset{z2 = z2 \cdot i}{\rightsquigarrow }
    \begin{pmatrix}
    -1 & -3i \\ 0 & 0 
    \end{pmatrix} \Rightarrow x_2 beliebig \\ \\
     x_1 & = -3i x_2 \\
      \mathcal{C} = \Biggl \{ \binom{-3i}{1} \Biggr \}
\end{align*}
\end{document}